\documentclass{article}
\usepackage{amsmath}

\title{A Study on Transformer Architectures for Natural Language Understanding}
\author{Alice Johnson and Bob Williams and Carol Davis}

\begin{document}
\maketitle

\begin{abstract}
This paper presents a comprehensive evaluation of transformer architectures for natural language understanding tasks. We compare BERT, GPT, and T5 variants across multiple benchmarks, demonstrating that task-specific fine-tuning strategies significantly impact performance.
\end{abstract}

\section{Introduction}
Natural language processing has undergone a paradigm shift with the introduction of transformer-based models. These architectures, initially designed for machine translation, have proven remarkably effective across a wide range of language understanding tasks.

\section{Methods}
We conducted experiments on four standard benchmarks: GLUE, SuperGLUE, SQuAD, and RACE. Each model was fine-tuned using three different strategies: full fine-tuning, adapter layers, and prompt tuning. All experiments were repeated five times with different random seeds.

\section{Results}
Our experiments show that T5-XXL achieves the highest average score across all benchmarks (91.3\%), followed by BERT-Large (87.6\%) and GPT-3 (85.2\%). Adapter-based fine-tuning consistently outperformed prompt tuning by 2-4 percentage points.

\section{Discussion}
The superior performance of T5 can be attributed to its encoder-decoder architecture, which enables more flexible sequence-to-sequence modeling. However, the computational cost of T5-XXL remains a significant barrier for practical deployment.

\section{Conclusion}
Transformer architectures continue to push the boundaries of NLP performance. Our findings suggest that the choice of fine-tuning strategy is as important as model selection for achieving optimal results.

\section{References}
\bibliographystyle{plain}
\bibliography{references}
\end{document}
